% Randomization locations in tree Bagging, random forests, and Extra Trees.
% Schematic only; no empirical quantities are encoded.
% Canvas: 164 mm x 69 mm. Dependencies: project palette and TikZ.
\begin{tikzpicture}[
  x=1mm,y=1mm,
  font=\footnotesize,
  text=BookInk,
  stage/.style={
    rectangle,
    draw=BookMuted,
    fill=BookPaper,
    line width=0.6pt,
    minimum width=39mm,
    minimum height=9mm,
    text width=35mm,
    align=center,
    inner sep=1mm
  },
  randomized/.style={
    stage,
    draw=FigureInput,
    fill=FigureInput!10,
    line width=0.9pt
  },
  searched/.style={
    stage,
    draw=FigureModel,
    fill=FigureModel!8,
    line width=0.8pt
  },
  output/.style={
    stage,
    draw=FigureOutput,
    fill=FigureOutput!6,
    line width=0.8pt
  },
  flow/.style={
    -{Latex[length=2.2mm,width=1.5mm]},
    draw=BookMuted,
    line width=0.8pt
  }
]
  \path[use as bounding box] (0,0) rectangle (164,69);

  \node[font=\heiti\small] at (27,65) {(a) 树Bagging};
  \node[randomized] (bag-sample) at (27,51) {随机化训练样本\\自助采样};
  \node[searched] (bag-feature) at (27,35) {检查全部候选特征};
  \node[searched] (bag-threshold) at (27,19) {搜索候选阈值};
  \node[output] (bag-tree) at (27,5) {得到一棵树};
  \draw[flow] (bag-sample)--(bag-feature);
  \draw[flow] (bag-feature)--(bag-threshold);
  \draw[flow] (bag-threshold)--(bag-tree);

  \node[font=\heiti\small] at (82,65) {(b) 随机森林};
  \node[randomized] (rf-sample) at (82,51) {随机化训练样本\\通常为自助采样};
  \node[randomized] (rf-feature) at (82,35) {随机化候选特征\\每个节点重新抽取};
  \node[searched] (rf-threshold) at (82,19) {在特征子集内搜索阈值};
  \node[output] (rf-tree) at (82,5) {得到一棵树};
  \draw[flow] (rf-sample)--(rf-feature);
  \draw[flow] (rf-feature)--(rf-threshold);
  \draw[flow] (rf-threshold)--(rf-tree);

  \node[font=\heiti\small] at (137,65) {(c) Extra Trees};
  \node[stage] (extra-sample) at (137,51)
    {训练样本\\原始方法使用完整训练集};
  \node[randomized] (extra-feature) at (137,35)
    {随机化候选特征};
  \node[randomized] (extra-threshold) at (137,19)
    {随机生成候选阈值\\再按准则择优};
  \node[output] (extra-tree) at (137,5) {得到一棵树};
  \draw[flow] (extra-sample)--(extra-feature);
  \draw[flow] (extra-feature)--(extra-threshold);
  \draw[flow] (extra-threshold)--(extra-tree);

  \draw[BookRule,line width=0.6pt] (54.5,1)--(54.5,66);
  \draw[BookRule,line width=0.6pt] (109.5,1)--(109.5,66);
\end{tikzpicture}
